

\section{Le poids des mots}

\begin{frame}{Pondération des termes}

Le classement s'appuie sur l'idée qu'il est possible 
d'identifier \alert{l'importance des termes} dans un document. Deux
idées essentielles:

\vfill

Plus un terme est \alert{fréquent dans un document}, plus il est représentatif du contenu du document.

\begin{itemize}
   \item Ex.: \emph{contrepoint} apparaît 10 fois dans le document $d$, $\Rightarrow$ $d$
      est un document important pour une recherche sur ``contrepoint''.
      \item \alert{Normalisation}\,: si $d$ est très long, il est normal de trouver plus d'occurrences
         des termes\,; on peut décider de \emph{normaliser} cet indicateur pour éviter ce biais.
\end{itemize}

\vfill

Plus un terme est \alert{rare dans la collection}, plus sa présence dans un document est importante.

\begin{itemize}
   \item Ex.: \emph{musicologie} apparaît 100 fois dans une grande collection, 4 fois dans le document $d$ $\Rightarrow$ $d$
      est un document important pour une recherche sur ``musicologie''.
\end{itemize}

\end{frame}


\begin{frame}[fragile]{Premier indicateur\,: la \alert{fréquence du terme}}

Soit un document $d$, $t$ un terme. 

\vfill

La \alert{fréquence} du terme $t$ dans $d$,  noté $n_{t,d}$, est simplement
le nombre d'occurrences de $t$ dans $d$.

\begin{noteblock}{Attention}
On parle bien de \alert{termes}, résultats de la normalisation lexicale (racinisation,
élimination des mots inutiles, etc.)
\end{noteblock}

\vfill

Exemple\,: soit $d$ le document  suivant\,:

\begin{smallverbatim}
Spider Cochon Spider Cochon, il peut marcher au plafond, 
Est ce qu'il peut faire une toile ? Bien sûr que non, 
c'est un cochon. Prends garde ! Spider Cochon est là !
\end{smallverbatim}

\medskip

Donc $tf(\rm{cochon}, d)=4$

\end{frame}


\begin{frame}{Deuxième indicateur\,: \alert{fréquence inverse des documents}}

Soit $t$ un terme, une collection $D$. On mesure sa \alert{rareté} de $t$
par l'inverse de sa fréquence dans $D$.

\begin{itemize}

\item Le nombre total de documents est $|D|$ 

\item le nombre de documents avec $t$ est $|\{d'\in D, |, n_{t,d'}>0\}|$


\end{itemize}

\vfill

La rareté de $t$ est donc mesurée par\,:

$$\frac{|D|}{|\{d'\in D\,|\,n_{t,d'}>0\}|}$$

\vfill

\alert{Ajustement}. La valeur obtenu par la formule ci-dessus croît très vite avec la
taille de la collection. On ajuste en prenant le logarithme et on obtient
la \alert{fréquence inverse des documents} (\emph{inverse document frequency}, idf) 

\[\mathop{\mathrm{idf}}(t)=\log\frac{|D|}{\left|\left\{d'\in
D\,|\,n_{t,d'}>0\right\}\right|}.\]

\end{frame}

\begin{frame}{Pondération par TF-IDF}

Le poids d'un terme dans un document est représenté par l'indicateur  \alert{Term Frequency---Inverse Document Frequency} (tf-idf)
 \[
\mathop{\mathsf{tf{}idf}}(t,d)={\color{magenta}n_{t,d}}\cdot
\color{orange}\log\frac{|D|}{\left|\left\{d'\in D\,|\,n_{t,d'}>0\right\}\right|}\]
\begin{center}\begin{tabular}{cl}
$n_{t,d}$& nombre d'occurrences de $t$ dans $d$\\
$D$& ensemble de tous les documents
\end{tabular}
\end{center}

\begin{itemize}
  \item le tf-idf décroît quand un terme est présent dans beaucoup de documents\,;
  \item il décroît également quand il est peu présent dans un document\,;
  \item il est maximal pour les termes peu fréquents apparaissant beaucoup dans un document particulier.
\end{itemize}

\begin{block}{Descripteurs}
Le tf.idf remplace l'indicateur 0/1 dans la matrice d'incidence.
\end{block}
\end{frame}


\begin{frame}[fragile]{Pondérons nos cochons}

Voici des documents A, B et C.

\begin{smallverbatim}
Spider Cochon Spider Cochon, il peut marcher au plafond, 
Est ce qu'il peut faire une toile ? Bien sûr que non, 
c'est un cochon. Prends garde ! Spider Cochon est là !
\end{smallverbatim}

\begin{smallverbatim}
Un petit cochon, pendu au plafond
\end{smallverbatim}

\begin{smallverbatim}
Les Trois Petits Cochons est un conte traditionnel européen 
mettant en scène trois jeunes cochons et un loup.
\end{smallverbatim}

\begin{exoblock}
Calculer le tf et l'idf pour les termes "cochon", "loup" et "plafond" 
et pour chaque document. En déduire la matrice d'incidence.
\end{exoblock}

% Idf(cochon)=1; idf(plafond) = 3/2, idf(loup)=3
% Pour A: cochon 4 | plafond 1 * 3/2 | loup 0
% Pour B: cochon 1 | plafond 1 *  3/2 | loup 0
% Pour C: cochon 2 | plafond 0 | loup 1 * 4
\end{frame}


\section{Similarité basée sur le tf-idf}


\begin{frame}{Calcul de la similarité}

À ce stade, on peut \alert{décrire} un document par un \alert{vecteur} composé
des \alert{poids} de tous ses termes. 

\medskip

\begin{center}
Descr(C) = ['cochon': 2, 'plafond': 0, 'loup': 4]
\end{center}

\vfill

Le poids est à 0 pour les termes qui n'apparaissent pas dans le document. On a un 
espace vectoriel $E = \mathbb{N}^{|V|}$, $V$ étant le vocabulaire.

\medskip
\begin{center}
['cochon': 0,  ..., 'jaguar': 4, ..., 'loup': 0, ..., 'python': 2, ...., 'mouton': 0]
\end{center}

\begin{itemize}
  \item L'espace a \alert{beaucoup} de dimensions (des millions d'axes / termes). 
   \item Chaque vecteur 
est principalement constitué de 0.
\end{itemize}

\begin{noteblock}{Distance Euclidienne?}
Potentiellement très coûteuse à calculer, et introduit un biais lié à la longueur
des documents. 
\end{noteblock}
\end{frame}


%%%%%%%%
\begin{frame}{Classement par cosinus}

Plus deux documents sont ``proches'' l'un de l'autre, plus l'angle de leurs vecteurs descripteurs est petit. 

\vfill

 
\figSlideWithSize{CosineSim}{7cm}



\alert{Rappel}\,: le cosinus est une fonction décroisante sur l'intervalle $[0, 90]$.

\end{frame}


\begin{frame}{La similarité cosinus}

La \alert{similarité cosinus} est un bon candidat pour mesurer la proximité des
vecteurs dans $E\mathbb{N}^{|V|}$


\vfill

\begin{itemize}
  \item Indifférent la la \alert{longueur} (norme) des vecteurs.
  \item Maximal si même direction (angle = 0, cosinus = 1)
  \item Minimal si directions "orthogonales" (pas de terme en commun)
  \item Varie continuement entre 0 et 1.
  
\end{itemize}

\begin{block}{En pratique}
Calcul efficace car nécessite seulement les coordonnées non nulles.
\end{block}
\end{frame}


\begin{frame}{Calcul du cosinus de deux vecteurs}

Produit scalaire de deux vecteurs (niveau 3ème !)

\[
    v_1 . v_2  = ||v_1|| \times ||v_2|| \times cos \theta = \sum_{i=1}^n v_1[i] \times v_2[i]
\]
où $\theta$ désigne l'angle entre les deux vecteurs et $||v||$ la norme d'un vecteur.

\vfill

Donc:

\[
   cos { } \theta  = \frac{\sum_{i=1}^n v_1[i] \times v_2[i]}{||v_1|| \times ||v_2||}
\]

\begin{block}{Normalisation}
La division par la norme revient à éliminer le biais lié à la longueur des documents.
\end{block}
\end{frame}


\begin{frame}{Calcul du cosinus, en pratique}

\alert{Première étape}\,: on calcule la \alert{norme} du vecteur 
représentant chaque document, on la stocke.
\medskip

Norme du vecteur $\vec{d}$:

      $$||\vec{d}|| = \sqrt{\Sigma_i d_i^2}$$
      
\vfill
      
\alert{Seconde étape}\,: $d$ le document, $q$ la requête; on prend
les vecteurs $\vec{q}$ (calculé à la volée) et $\vec{d}$ (stocké dans un \alert{index}). 
      
\[
      \frac{1}{||\vec{d}||} \times \frac{1}{||\vec{q}||}         \times \Sigma_i d_i q_i
\]      
      
\begin{block}{Approximation}
On peut ignorer la norme de la requête (qui est constante), et la racine carré
pour la calcul des normes: c'est le classement qui nous intéresse.
\end{block}
           
\end{frame}

